\chapter{Pretraining on SpectralEarth}\label{ch:pretrain}
\begingroup\setlength{\emergencystretch}{5em}% long provenance paths: allow loose lines rather than overfull ones
\def\tbrk{\hfil\penalty50\hfilneg}% local: a line may break before a provenance tag, leaving the line ragged

\begin{onething}
A long pretraining run is only as good as three quiet decisions: the number that ranks its epochs must be
measured on the same scale every epoch; a stem that \emph{can} read any band list must be \emph{trained} on many
band lists, or it reads only the one it saw; and data that does not fit on the disk is walked through in blocks
on one unbroken learning-rate schedule. This chapter derives each decision from numbers, then follows the
pretrained model onto the three sensors it will be fine-tuned on.
\end{onething}

\section{What this chapter is about}\label{sec:pretrain:intro}

\Cref{ch:trainloop} built a training loop: batches of patches go through the model, a loss is computed, the
gradient of the loss nudges every parameter, and after each \emph{epoch} (one pass over the epoch's training
patches) the loop measures a \emph{validation loss} on patches it never trains on. \Cref{ch:losses} built the loss
itself, $L1 + \lambda(t)\,\SAM$, with the weight $\lambda$ of the spectral angle \citep{kruse1993spectral} ramping up from 0 over the first
epochs. \Cref{ch:specdecoder} built the reduced-attention spectral stem, which reads a pixel's bands as tokens
that carry their own wavelengths, so that in principle it can read any sensor's band list.

This chapter takes that model and trains it for real, on \emph{SpectralEarth}\index{SpectralEarth}: about half a
million hyperspectral patches that the German space agency DLR cut from EnMAP scenes and published for exactly
this purpose \citep{braham2024spectralearth}. Three things went wrong, or nearly wrong, on the way, and each
teaches a piece of mathematics:
\begin{enumerate}
  \item \textbf{The ruler moved.} The validation loss was computed with the training loss's ramped weight. A score
    whose formula changes from epoch to epoch cannot rank epochs (\cref{sec:pretrain:ruler}).
  \item \textbf{One band list is not any band list.} A stem trained on SpectralEarth's single list of 202 bands read
    PRISMA's own 185 bands badly, worst over water. The cause is an average that gives every band one vote, and a
    set of small wavelength networks that were only ever evaluated at one list (\cref{sec:pretrain:density}).
  \item \textbf{The data did not fit.} The final run, F001, needed ten times more patches than the disk could hold.
    It walked through them in blocks of 30,000 patches, with arithmetic that decides how much it sees, how much it
    stores and when it rotates (\cref{sec:pretrain:f001}).
\end{enumerate}
The chapter ends with the three real sensors that fine-tuning (\cref{ch:finetune}) will read at their own bands:
EnMAP, PRISMA and AVIRIS-NG, each with quirks that one rule handles (\cref{sec:pretrain:native}).

\section{SpectralEarth in numbers}\label{sec:pretrain:se}

Before training anything, count what you have. The data layer lives in \code{app/utils/spectralearth/} and is
documented in \file{docs/lld/spectralearth.md}.

\paragraph{Patches.} SpectralEarth holds 539,755\tbrk\tagM{docs/lld/spectralearth.md} patches. Each
\index{patch}patch is a small square cut from a scene: 128 $\times$ 128 pixels, each pixel a spectrum of 202 bands.
The values are stored as 16-bit integers (\emph{int16}, 2 bytes each): reflectance multiplied by 10,000, with
$-32768$ meaning ``no data''. So reflectance $0.1934$ is stored as $1934$.

\paragraph{Why 202 bands.} EnMAP records 224 bands. SpectralEarth drops two runs of them, the bands where water
vapour in the air absorbs almost all sunlight so that the ground is invisible: bands 127 to 141 and 161 to 167
(1-based).\tbrk\tagM{docs/lld/spectralearth.md} Count each run. A run from $a$ to $b$ inclusive has $b - a + 1$
members:
\begin{align*}
  \text{run 1:}\quad 141 - 127 + 1 &= 15 \text{ bands}\\
  \text{run 2:}\quad 167 - 161 + 1 &= 7 \text{ bands}\\
  \text{kept:}\quad 224 - 15 - 7 &= 202 \text{ bands.}
\end{align*}
The data layer derives this list from a table and refuses anything else; the LLD's first invariant is
``never hard-code a SpectralEarth wavelength; read the table''.

\begin{worked}[label={we:pretrain:patchbytes}]{The size of one patch}
Do this on paper. A patch is 202 bands of 128 $\times$ 128 pixels at 2 bytes per number.
\begin{align*}
  \text{pixels per band:}\quad 128 \times 128 &= 16{,}384\\
  \text{numbers per patch:}\quad 202 \times 16{,}384 &= 3{,}309{,}568\\
  \text{bytes per patch:}\quad 3{,}309{,}568 \times 2 &= 6{,}619{,}136 \text{ bytes}\\
  \text{in megabytes (}10^6\text{):}\quad 6{,}619{,}136 / 1{,}000{,}000 &\approx 6.62 \text{ MB.}
\end{align*}
A \emph{shard}\index{shard} is a file that bundles 150 patches (a tar archive, read in order). One shard is therefore
\begin{align*}
  150 \times 6{,}619{,}136 &= 992{,}870{,}400 \text{ bytes} \approx 0.99 \text{ GB,}
\end{align*}
which matches the LLD's ``150-patch shards (\textasciitilde 0.99 GB)''.\tbrk\tagM{docs/lld/spectralearth.md}
\end{worked}

\paragraph{The split.} Patches from one place on two dates look alike. If one date were in training and the other in
validation, the validation score would flatter the model. So the split is \emph{by site}: a site goes to the test
side when a hash of its name, read as a number between 0 and 1, falls below the test fraction $f$. Every date of a
site lands on the same side, and the same seed gives the same split on any machine. The \code{se-v1} split at
$f = 0.05$ gave 512,849 training and 26,906 test patches.\tbrk\tagM{docs/lld/spectralearth.md} Check that it adds up
and that the test share is close to 5\%:
\begin{align*}
  512{,}849 + 26{,}906 &= 539{,}755 \quad\text{(every patch, once)}\\
  26{,}906 / 539{,}755 &= 0.0498 \approx 5\%.
\end{align*}
The share is not exactly 5\% because sites hold different numbers of patches; the hash chooses sites, not patches.

\paragraph{The normalisation curve.} The sharder also writes \code{stats.json}: for every one of the 202
wavelengths, the mean and standard deviation of valid reflectance over the training split. This is the
\emph{normalisation curve}\index{normalisation curve} of \cref{ch:bandpipeline}: inside the model each band value $v$ at wavelength $\lambda$
becomes $(v - \mu(\lambda)) / \sigma(\lambda)$, ``standard deviations from that band's mean''. Keep this curve in
mind; \cref{ch:finetune} shows why it must never be recomputed after pretraining.

\section{A ruler that does not move}\label{sec:pretrain:ruler}

\subsection{Two jobs for one number}

During pretraining the validation loss has a second job besides reporting progress. It decides two things:
\begin{itemize}
  \item \textbf{\texttt{best.pt}}\index{best checkpoint}: which epoch's weights are kept as ``the best''. A
    \emph{checkpoint} is a file holding the model's weights (and the optimiser's state) at the end of an epoch;
    \code{best.pt} is the one with the lowest validation score so far.
  \item \textbf{Early stopping}: whether to stop the run before its last epoch.
\end{itemize}

\begin{definition}[Early stopping]\index{early stopping}\index{patience}
Keep the best validation score seen so far, $S^\ast$, and a counter of \emph{bad epochs}. After each epoch with
score $S_e$: if $S_e < S^\ast(1 - \delta)$, the epoch is an improvement, so set $S^\ast \gets S_e$ and reset the
counter to 0; otherwise add 1 to the counter. Stop when the counter reaches the \emph{patience} $P$, but never before
a minimum epoch $E_{\min}$. The number $\delta$ is the \emph{minimum relative improvement}.
\end{definition}

The rule in \code{app/foundation_models/trainers/spectral_stem_mae_trainer.py} is exactly this:
\begin{lstlisting}
if self._best_val == math.inf or avg_val < self._best_val * (1 - cfg.early_stop_min_delta):
    self._best_val, self._bad_epochs = avg_val, 0   # improvement: new best, count reset
else:
    self._bad_epochs += 1                           # not better by delta: one more bad epoch
self._stop = bool(cfg.early_stop_patience) and epoch >= cfg.early_stop_min_epochs \
    and self._bad_epochs >= cfg.early_stop_patience
\end{lstlisting}
The experiment configs (\file{research/experiments/indradhanu-se-pretrain/E00*/config.json}) set $\delta = 0.01$,
$P = 10$, $E_{\min} = 20$, a final SAM weight of 0.5 and a ramp of 10 epochs.\tbrk\tagM{research/experiments/indradhanu-se-pretrain/F001-final-pretrain/config.json}
With $\delta = 0.01$ a score counts as an improvement only if it is below $0.99\,S^\ast$: one percent better than the
best.

\subsection{The trap, as algebra}

The training loss of \cref{ch:losses} is
\begin{equation}\label{eq:pretrain:trainloss}
  \mathcal{L}_e = L1_e + \lambda_e\,\SAM_e, \qquad \lambda_e = \lambda_{\max}\min\!\left(1, \frac{e-1}{R}\right),
\end{equation}
where $e = 1, 2, \dots$ is the epoch as a person counts it, $\lambda_{\max} = 0.5$ and $R = 10$. The trainer counts
epochs from 0 internally (\code{_get_sam_weight(epoch)} returns $\lambda_{\max}\min(1, \texttt{epoch}/R)$ with
$\texttt{epoch} = e - 1$), which is why $e - 1$ appears. Work four values:
\begin{align*}
  e = 1:&\quad 0.5 \times \min(1, 0/10) = 0.5 \times 0.0 = 0.00\\
  e = 3:&\quad 0.5 \times \min(1, 2/10) = 0.5 \times 0.2 = 0.10\\
  e = 6:&\quad 0.5 \times \min(1, 5/10) = 0.5 \times 0.5 = 0.25\\
  e = 11:&\quad 0.5 \times \min(1, 10/10) = 0.5 \times 1.0 = 0.50.
\end{align*}

The first validation code scored epochs with the same formula, ramp included. Ask when that score can rise from one
epoch to the next even though both measurements fall. Write the change in score:
\begin{align*}
  S_{e+1} - S_e
  &= \bigl(L1_{e+1} + \lambda_{e+1}\SAM_{e+1}\bigr) - \bigl(L1_e + \lambda_e\SAM_e\bigr).
\intertext{Add and subtract $\lambda_e \SAM_{e+1}$, which changes nothing, so that each difference appears on its
own:}
  &= (L1_{e+1} - L1_e) + \lambda_e(\SAM_{e+1} - \SAM_e) + (\lambda_{e+1} - \lambda_e)\SAM_{e+1}.
\intertext{Name the three differences $\Delta L1$, $\Delta\SAM$ and $\Delta\lambda$:}
  S_{e+1} - S_e &= \underbrace{\Delta L1}_{\le 0\text{ if L1 improves}}
     + \underbrace{\lambda_e\,\Delta\SAM}_{\le 0\text{ if SAM improves}}
     + \underbrace{\Delta\lambda\;\SAM_{e+1}}_{> 0\text{ while the ramp rises}}.
\end{align*}
\begin{keypoint}
While $\lambda$ is still rising, the third term is positive and grows with the size of SAM; if it outweighs the
first two, a model that improved on \emph{both} measurements scores \emph{worse}.
\end{keypoint}
SAM is measured in radians and is several times larger than L1 in reflectance (a SAM of 0.2 rad against an L1 of
0.05), so the third term easily wins.

\begin{worked}[label={we:pretrain:ramp}]{The same model scored two ways}
These values are illustrative \tagA, chosen to match what the live run E003 showed: a first-epoch score of about
0.055 and every later epoch about 0.1. L1 and SAM both improve every epoch.
\begin{center}
\begin{tabular}{rccc}
\toprule
epoch $e$ & L1 & SAM (rad) & $\lambda_e$\\
\midrule
1 & 0.055 & 0.20 & 0.00\\
3 & 0.045 & 0.17 & 0.10\\
6 & 0.040 & 0.15 & 0.25\\
11 & 0.034 & 0.13 & 0.50\\
\bottomrule
\end{tabular}
\end{center}
\textbf{Scored with the ramp}, $S_e = L1_e + \lambda_e\SAM_e$:
\begin{align*}
  e = 1:&\quad 0.055 + 0.00 \times 0.20 = 0.055 + 0.000 = 0.055\\
  e = 3:&\quad 0.045 + 0.10 \times 0.17 = 0.045 + 0.017 = 0.062\\
  e = 6:&\quad 0.040 + 0.25 \times 0.15 = 0.040 + 0.0375 = 0.0775\\
  e = 11:&\quad 0.034 + 0.50 \times 0.13 = 0.034 + 0.065 = 0.099.
\end{align*}
\textbf{Scored at the full weight}, $S_e = L1_e + 0.5\,\SAM_e$ every epoch:
\begin{align*}
  e = 1:&\quad 0.055 + 0.5 \times 0.20 = 0.055 + 0.100 = 0.155\\
  e = 3:&\quad 0.045 + 0.5 \times 0.17 = 0.045 + 0.085 = 0.130\\
  e = 6:&\quad 0.040 + 0.5 \times 0.15 = 0.040 + 0.075 = 0.115\\
  e = 11:&\quad 0.034 + 0.5 \times 0.13 = 0.034 + 0.065 = 0.099.
\end{align*}
Now check the decomposition on the step from epoch 1 to epoch 3 with the ramp. The three terms are
\begin{align*}
  \Delta L1 &= 0.045 - 0.055 = -0.010\\
  \lambda_1\,\Delta\SAM &= 0.00 \times (0.17 - 0.20) = 0.00 \times (-0.03) = 0.000\\
  \Delta\lambda\,\SAM_3 &= (0.10 - 0.00) \times 0.17 = 0.017\\
  \text{sum} &= -0.010 + 0.000 + 0.017 = +0.007,
\end{align*}
and indeed $0.062 - 0.055 = 0.007$. The ramp term ($+0.017$) is larger than the whole improvement in L1
($-0.010$), so the score rises. The model is identical in both columns; only the ruler differs.
\end{worked}

\coursefig[0.85\linewidth]{iw-41-val-score}{The four epochs of \cref{we:pretrain:ramp} scored two ways. With the
training ramp (amber, dashed) the validation score climbs from 0.055 to 0.099 although L1 and SAM both fall; epoch 1
stays ``best'' for ever. At the full SAM weight (teal) the score falls from 0.155 to 0.099. The two meet at epoch 11,
where the ramp ends.}{fig:pretrain:valscore}

\subsection{What early stopping would have done}

Run the rule on the ramped scores. Epoch 1 sets $S^\ast = 0.055$, so a later epoch must score below
\[
  0.055 \times (1 - 0.01) = 0.055 \times 0.99 = 0.05445.
\]
Every later epoch scores near 0.1, so none qualifies:
\begin{center}
\begin{tabular}{rll}
\toprule
epoch & below 0.05445? & bad epochs, and the stop test\\
\midrule
2 & no & 1\\
3 & no & 2\\
$\vdots$ & & \\
11 & no & 10: patience reached, but $11 < 20$, so no stop\\
$\vdots$ & & \\
20 & no & 19: $20 \ge 20$ and $19 \ge 10$, \textbf{stop}\\
\bottomrule
\end{tabular}
\end{center}
The run would have stopped at epoch 20, the first epoch the rule allows, and kept epoch 1's weights as \code{best.pt},
whatever the model had learned. The minimum of 20 epochs delayed the damage but did not prevent it: the comparison
itself was broken.

\begin{analogy}
A golf club records each member's best round of the season. The first round is scored on putts only; from the second
round on, every stroke counts. The first round becomes everyone's ``best round'' for the rest of the year, however
much their game improves. Early stopping is the club secretary who, seeing no new best for ten rounds, cancels the
member's season.
\end{analogy}

\subsection{The fix, and making it safe to resume into}

One line in \code{compute_validation_loss}: training keeps the ramp, validation uses the final weight from the
first epoch.
\begin{lstlisting}
return l1 + cfg.sam_weight * sam, n          # validation: lambda = 0.5 at every epoch
\end{lstlisting}
The source records why in its own comment: with the ramp, ``epoch 1 (lambda 0) scored \textasciitilde 0.055 and every
later epoch \textasciitilde 0.1 - early stop would have fired at the first allowed epoch whatever the model did''
(found in E003 on 2026-10-14). Two smaller pieces make the fix safe:
\begin{itemize}
  \item \textbf{A checkpoint remembers how it was scored.} The early-stop state saved inside each checkpoint carries
    \code{"score": "l1+full_sam"} (the constant \code{VAL_SCORE}). On resume, a stored state without that tag is
    ignored rather than compared with numbers on another scale.
  \item \textbf{The parts are logged separately.} \code{metrics.csv} records \code{val_l1}, \code{val_sam_deg} and
    \code{sam_weight} each epoch, so you can recompute any score by hand. Convert the angle first:
    radians $=$ degrees $\times \pi / 180$, so $4.90^\circ = 4.90 \times 3.1416 / 180 = 0.0855$ rad.
\end{itemize}

\begin{sourcenote}[: the validation score and the resume check]
Verified in \file{app/foundation_models/trainers/spectral_stem_mae_trainer.py}: \code{VAL_SCORE = "l1+full_sam"}
(line 44); \code{compute_validation_loss} returns \code{l1 + cfg.sam_weight * sam} (line 341) while
\code{compute_loss} returns \code{l1 + self._get_sam_weight(self._current_epoch) * sam} (line 315); a loaded
early-stop state whose \code{score} differs from \code{VAL_SCORE} is discarded with a log line (lines 148--151). The
ramp itself is \code{sam_max * min(1.0, epoch / ramp_epochs)} with a 0-based epoch
(\file{hyperspectral_segformer_mae_trainer.py}, \code{_get_sam_weight}).
\end{sourcenote}

\section{One band list is not any band list}\label{sec:pretrain:density}

\subsection{The symptom}

E004, the reduced-attention stem of \cref{ch:specdecoder}, was scored on the four labelled PRISMA scenes. A
\emph{ROC-AUC} (\cref{ch:measures}) of 0.5 means ``no better than a coin''. Scene set1 scored 0.217: worse than a coin,
so the model ranked real anomalies \emph{below} ordinary background.\tbrk\tagM{indradhanu workings/42-one-band-list.md}
Set1 is 75\% sea. Splitting its mean reconstruction error by surface showed where the problem sat:
\begin{center}
\begin{tabular}{lcc}
\toprule
& mean L1 on water & on land\\
\midrule
E004, PRISMA's native bands & 0.0462 & 0.0282\\
v0.2.0 (production) & 0.0046 & 0.0113\\
\bottomrule
\end{tabular}
\end{center}
E004's error on water was $0.0462 / 0.0046 \approx 10$ times production's, and 93\% of its highest-scoring pixels were
sea.

\paragraph{The cheap explanation, measured first.} ``SpectralEarth is land, so the model never learned water.'' A census
counted water: 6.8\% of EnMAP training tiles had more than 30\% water, against 6--12\% of a 3,000-patch SpectralEarth
sample.\tbrk\tagM{research/experiments/indradhanu-se-pretrain/E004-reduced-attention/} SpectralEarth has as much water as
EnMAP. Measuring the cheap explanation before building its fix saved a training run that would have changed nothing.

\paragraph{The clue.} The PRISMA bench scores each scene two ways: on PRISMA's \emph{native} bands (its own list, 185
bands after the product's default band filter) and \emph{resampled} onto SpectralEarth's 202 wavelengths. At epoch 18,
set1 scored 0.205 native and 0.725 resampled. Same scene, same model; only the band list changed.

\subsection{One vote per band}

\Cref{ch:specdecoder} gave the reduced stem two routes from a spectrum to its 32 channels, added together. Route 1
averages the bands inside each 40\,nm \emph{bin} and attends across the bins. Route 2, the \emph{skip path}, sends every
band straight to the channels:
\begin{equation}\label{eq:pretrain:skip}
  \vect{s} = \frac{1}{N}\sum_{b=1}^{N} v_b\,\vect{W}(\lambda_b),
\end{equation}
where $v_b$ is the normalised value of band $b$, $\lambda_b$ its wavelength, and $\vect{W}(\lambda)$ a weight vector
produced by a small network from the wavelength. Because $\vect{W}$ is computed from $\lambda$, the formula works for any
list of bands. But look at the factor $1/N$: every band gets one equal vote.

To see what that does, set every weight to 1 (one channel, $\vect{W} = 1$). Then \eqref{eq:pretrain:skip} is the plain
mean of the band values. Suppose a spectrum takes value $a$ in the visible and $c$ elsewhere, and a fraction $f$ of the
sensor's bands lie in the visible:
\begin{align*}
  s &= \frac{1}{N}\Bigl(\underbrace{a + \dots + a}_{fN \text{ times}} + \underbrace{c + \dots + c}_{(1-f)N\text{ times}}\Bigr)
\intertext{Add the repeated terms:}
    &= \frac{1}{N}\bigl(fN a + (1 - f) N c\bigr)
\intertext{Cancel $N$:}
    &= f a + (1 - f) c.
\end{align*}
The skip path's output depends on $f$, the share of bands in the visible: a property of the \emph{sensor}, not of the
ground.

\begin{worked}[label={we:pretrain:votes}]{One spectrum, two sensors, two answers}
Do this on paper. The spectrum is the same for both sensors: value 2 in the visible, $-1$ in the shortwave infrared.
Sensor A has bands at 500, 505 and 2000\,nm; sensor B at 500 and 2000\,nm.
\begin{align*}
  \text{A:}\quad s &= \frac{2 + 2 + (-1)}{3} = \frac{3}{3} = 1.0\\
  \text{B:}\quad s &= \frac{2 + (-1)}{2} = \frac{1}{2} = 0.5.
\end{align*}
Check with the formula $s = fa + (1-f)c$. For A, $f = 2/3$: $\tfrac{2}{3}\times 2 + \tfrac{1}{3}\times(-1) =
\tfrac{4}{3} - \tfrac{1}{3} = 1.0$. For B, $f = 1/2$: $\tfrac{1}{2}\times 2 + \tfrac{1}{2}\times(-1) = 1 - 0.5 = 0.5$.
Route 1 does not have this problem: it averages \emph{inside} a bin first, so 500 and 505\,nm become one bin value of 2
for both sensors.
\end{worked}

The real lists differ by more than the toy:\tbrk\tagM{indradhanu workings/42-one-band-list.md}
\begin{center}
\begin{tabular}{lcccc}
\toprule
& bands & visible ($<$700\,nm) & NIR (700--1000) & SWIR ($>$1000)\\
\midrule
SpectralEarth (trained on) & 202 & 53 (26\%) & 48 (24\%) & 101 (50\%)\\
PRISMA native (set1) & 185 & 30 (16\%) & 23 (12\%) & 132 (71\%)\\
\bottomrule
\end{tabular}
\end{center}
On PRISMA's list the skip path hears the visible at about $16/26 = 0.62$ of the volume it was trained with. Water's
signal is almost all in the visible (beyond about 750\,nm water is nearly black), so water pixels are distorted most.

\coursefig[0.95\linewidth]{iw-42-band-density}{Left: how each sensor's bands divide between the visible, near-infrared
and shortwave infrared. PRISMA's native list gives the visible 16\% of its bands against SpectralEarth's 26\%. Right:
E004's error on water, on SpectralEarth test patches with all 202 bands, on the same pixels thinned to PRISMA's
spacing, and with the skip path reweighted at inference.}{fig:pretrain:density}

\subsection{Proving it one variable at a time}

Two experiments isolated the band list from everything else, on E004's epoch-18 checkpoint (2026-10-14):
\begin{itemize}
  \item \textbf{Thinning.} Take SpectralEarth test patches (the model's own sensor and wavelengths) and keep only the 151
    of 202 bands nearest PRISMA's. No new sensor, no new wavelength: only fewer bands, spaced like PRISMA.
  \item \textbf{Reweighting.} At inference, scale each band's vote in the skip path by $n_{\text{train}}(\text{bin}) /
    n_{\text{now}}(\text{bin})$. A bin that had 8 bands in training and has 5 now gives each of its 5 bands weight
    $8/5 = 1.6$.
\end{itemize}
\begin{center}
\begin{tabular}{lcc}
\toprule
condition (mean L1) & water & land\\
\midrule
SpectralEarth, all 202 bands & 0.0155 & 0.0157\\
same pixels, thinned to PRISMA spacing & 0.0309 & 0.0226\\
thinned $+$ skip reweighted & 0.0203 & 0.0205\\
set1 native & 0.0462 & 0.0282\\
set1 native $+$ skip reweighted & 0.0352 & 0.0200\\
set1 resampled to SpectralEarth's wavelengths & 0.0138 & 0.0160\\
\bottomrule
\end{tabular}
\end{center}
Source:\tbrk\tagM{research/experiments/indradhanu-se-pretrain/E004-reduced-attention/} (as quoted in IW42).

Read it in order. Thinning doubles the water error: $0.0309 / 0.0155 = 1.99$. Reweighting closes a fraction of that gap,
\begin{align*}
  \text{gap} &= 0.0309 - 0.0155 = 0.0154\\
  \text{closed} &= 0.0309 - 0.0203 = 0.0106\\
  \text{fraction} &= 0.0106 / 0.0154 = 0.69,
\end{align*}
about two thirds, so the one-vote-per-band average is a real cause. It is not the only one: no fix brings set1 down to
the resampled 0.0138.

\subsection{The root cause and the three fixes}

Every wavelength-dependent piece of the stem (the in-bin kernel $k(\lambda)$, the skip weights $\vect{W}(\lambda)$, the
decoder's $\vect{U}(\lambda)$) is a small network fed the band's wavelength. In pretraining those networks were only ever
evaluated at one list: the same 202 wavelengths, every batch, for 18 epochs. ``Reads any band list'' was a property of the
\emph{architecture}, never of the \emph{training}.
\begin{keypoint}
A function learned from examples is only trustworthy near the examples. A network of the wavelength that has only seen
one list of 202 wavelengths has learned nothing about what a missing neighbour means.
\end{keypoint}

The fixes adopted on 2026-10-14:

\paragraph{Fix 1: train on many band lists.}\index{band-list augmentation} With probability \code{band_augment_prob}
(0.8 in F001) a batch reads a random subset of its bands, in one of three shapes chosen with equal probability
(\code{sample_band_subset} in \file{spectral_stem_mae_trainer.py}):
\begin{enumerate}
  \item \textbf{thinning}: keep each band with a probability drawn uniformly between 0.5 and 0.95;
  \item \textbf{gaps}: remove one to three windows, each 20--150\,nm wide, at random centres;
  \item \textbf{even grid}: keep the band nearest each point of a grid with a random step of 6--15\,nm (PRISMA-like at
    10\,nm, EnMAP-like at 6.5\,nm).
\end{enumerate}
If fewer than 40 bands survive, the batch keeps all of them. The kept wavelengths are also jittered by up to
$\pm 3$\,nm (\code{band_jitter_nm}). The model still reconstructs \emph{all} the bands, so it learns to fill in the ones
it was not given. Work the expected size of a thinned list on SpectralEarth: the keep probability averages
$(0.5 + 0.95)/2 = 0.725$, so on average $202 \times 0.725 = 146.45$, about 146 bands are read. An even grid of 10\,nm over
SpectralEarth's range (about 418 to 2445\,nm, \tagA\ rounded) has about $(2445 - 418)/10 = 202.7$ points, but the bands
inside the dropped water-vapour regions are missing, so several grid points share a nearest band and the kept list is
shorter than the grid.

\paragraph{Fix 2: one vote per nanometre.}\index{band width weighting} Weight each band by the width of spectrum it stands
for, $w_b$, and divide by the total width:
\begin{equation}\label{eq:pretrain:width}
  s_w = \frac{\sum_b w_b\, v_b}{\sum_b w_b}.
\end{equation}
Why this is the right average: if the spectrum is a function $v(\lambda)$ and band $b$ covers an interval of width $w_b$
around $\lambda_b$, then $w_b v_b$ approximates the area under $v$ over that interval. The sum is a \emph{Riemann sum}\index{Riemann sum}, an
approximation of the integral:
\begin{align*}
  \sum_b w_b v_b &\approx \int_{\lambda_{\min}}^{\lambda_{\max}} v(\lambda)\,d\lambda,
  \qquad \sum_b w_b \approx \lambda_{\max} - \lambda_{\min},
\intertext{so}
  s_w &\approx \frac{1}{\lambda_{\max} - \lambda_{\min}}\int_{\lambda_{\min}}^{\lambda_{\max}} v(\lambda)\,d\lambda,
\end{align*}
the mean value of the spectrum over wavelength. That quantity belongs to the ground, not to the sensor: two sensors that
cover the same range at different densities approximate the same integral.

\begin{worked}[label={we:pretrain:widths}]{The toy again, with widths}
Use idealised widths \tagA: in A each visible band covers 5\,nm; in B the one visible band covers 10\,nm; each infrared
band covers 10\,nm.
\begin{align*}
  \text{A:}\quad s_w &= \frac{2\times 5 + 2 \times 5 + (-1)\times 10}{5 + 5 + 10}
                     = \frac{10 + 10 - 10}{20} = \frac{10}{20} = 0.5\\
  \text{B:}\quad s_w &= \frac{2 \times 10 + (-1)\times 10}{10 + 10} = \frac{20 - 10}{20} = \frac{10}{20} = 0.5.
\end{align*}
Same spectrum, same answer, and no new parameters.
\end{worked}

\begin{sourcenote}[: how the real widths are computed]
The course's toy uses idealised widths. The code (\code{band_widths} in
\file{app/foundation_models/components/spectral_reduced.py}, lines 39--53) sorts the wavelengths, takes the gaps between
neighbours, \emph{caps each gap at 20\,nm} (\code{MAX_GAP_NM}, ``a band beside an excluded region would otherwise claim the
whole gap''), gives an interior band the mean of its two gaps, and gives each end band its one gap in full. On the toy's
three-band list A (500, 505, 2000\,nm) that gives gaps $5$ and $\min(1495, 20) = 20$, so widths $5$, $(5+20)/2 = 12.5$ and
$20$, and
\[
  s_w = \frac{2\times 5 + 2\times 12.5 + (-1)\times 20}{5 + 12.5 + 20} = \frac{10 + 25 - 20}{37.5} = \frac{15}{37.5} = 0.4,
\]
while B (gap capped at 20, widths 20 and 20) gives $(40 - 20)/40 = 0.5$. On a three-band toy the end rule and the cap
dominate; on real lists of 185--372 bands only the two end bands and the bands beside the water-vapour gaps are affected,
and the capped widths are the point: a band next to a 200\,nm hole must not claim to represent it.
\end{sourcenote}

\paragraph{Fix 3: ignore wavelengths the model never trained on.} The checkpoint carries its training wavelengths, so a
band farther than 15\,nm from all of them is dropped when a scene is loaded (\code{trained_band_mask}). The rule comes from
the checkpoint, not from a fixed grid. On the bench's 185 PRISMA bands, 165 pass and 20 are dropped: those inside
SpectralEarth's water-vapour gaps (1306--1502 and 1727--1959\,nm) and past its last band. \Cref{ch:finetune} shows how
fine-tuning grows that record.

\paragraph{Did it work?} E005 carried the three fixes. Extended to 20 epochs it read PRISMA's own bands about as well as E004
had read the resampled ones: mean ROC-AUC 0.752 against 0.754.\tbrk\tagM{research/experiments/indradhanu-se-pretrain/E005-band-robust/README.md}
That made E005's configuration the one to train properly.

\section{F001: two hundred epochs on rotating data}\label{sec:pretrain:f001}

\subsection{Why more data, not only more epochs}

F001 is the final SpectralEarth pretraining of the next Indradhanu, with E005's configuration exactly, and the
checkpoint all fine-tuning starts from. Its plan was 200 epochs. An epoch here is \code{train_samples_per_epoch} =
30,000 patches.\tbrk\tagM{research/experiments/indradhanu-se-pretrain/F001-final-pretrain/config.json}

The cheap way to train 200 epochs is to show the same 30,000 patches 200 times. A model with millions of parameters can
then \emph{memorise}\index{overfitting}: improve on those exact patches without improving on patches it has never seen.
F001 instead walks through ten \emph{blocks}\index{block (of shards)}. A block is 200 consecutive training shards; at 150
patches per shard that is $200 \times 150 = 30{,}000$ patches, exactly one epoch's worth. Block $k$ is shards $200k$ to
$200k + 199$ and trains epochs $20k + 1$ to $20(k+1)$:
\begin{center}
\begin{tabular}{rcc}
\toprule
block $k$ & shards & epochs\\
\midrule
0 & 0--199 & 1--20\\
1 & 200--399 & 21--40\\
2 & 400--599 & 41--60\\
$\vdots$ & $\vdots$ & $\vdots$\\
9 & 1800--1999 & 181--200\\
\bottomrule
\end{tabular}
\end{center}

\begin{worked}[label={we:pretrain:coverage}]{Coverage per epoch and per run}
Do this on paper. How often is each patch seen, how many different patches does the run see, and what share of the
training split is that?
\begin{align*}
  \text{epochs per block:}&\quad 200 / 10 = 20\\
  \text{times each patch is seen:}&\quad 20 \text{ (one per epoch of its block)}\\
  \text{distinct patches:}&\quad 10 \times 30{,}000 = 300{,}000\\
  \text{share of the split:}&\quad 300{,}000 / 512{,}849 = 0.585 \approx 58\%\\
  \text{patch presentations in total:}&\quad 200 \times 30{,}000 = 6{,}000{,}000.
\end{align*}
Compare the E-runs (E001--E005), capped at 20 epochs on one block: each patch also seen 20 times, but only 30,000 distinct
patches, $30{,}000/512{,}849 = 0.058$, under 6\% of the split. F001 repeats each patch exactly as often as the E-runs did,
and sees ten times as many of them.
\end{worked}

\subsection{Why not put all ten blocks on the disk}

Because they do not fit. From \cref{we:pretrain:patchbytes}, one shard is about 0.99\,GB, so
\begin{align*}
  \text{one block:}&\quad 200 \times 0.99 \text{ GB} = 198 \text{ GB}\\
  \text{all ten:}&\quad 10 \times 198 \text{ GB} = 1{,}980 \text{ GB} \approx 1.98 \text{ TB.}
\end{align*}
The L40S instance's local NVMe disk holds about 412\,GB. It fits two blocks, the one training and the one being downloaded
for next time: $2 \times 198 = 396$\,GB $< 412$\,GB. Three blocks ($594$\,GB) do not.

\subsection{One schedule across all ten blocks}

The learning rate follows the cosine schedule of \cref{ch:trainloop} \citep{loshchilov2017sgdr}, computed over all 200
epochs, with $\eta_0 = 0.0005$ and $\eta_{\min} = 10^{-6}$:
\begin{equation}\label{eq:pretrain:cosine}
  \eta(t) = \eta_{\min} + \tfrac{1}{2}(\eta_0 - \eta_{\min})\Bigl(1 + \cos\frac{\pi t}{200}\Bigr).
\end{equation}
Note $\eta_0 - \eta_{\min} = 0.0005 - 0.000001 = 0.000499$. At the first three rotations:
\begin{align*}
  t = 20:&\quad \cos(0.1\pi) = 0.9511\\
         &\quad \eta = 0.000001 + 0.5 \times 0.000499 \times 1.9511 = 0.000001 + 0.0004868 = 0.000488\\
  t = 40:&\quad \cos(0.2\pi) = 0.8090\\
         &\quad \eta = 0.000001 + 0.5 \times 0.000499 \times 1.8090 = 0.000001 + 0.0004513 = 0.000452\\
  t = 60:&\quad \cos(0.3\pi) = 0.5878\\
         &\quad \eta = 0.000001 + 0.5 \times 0.000499 \times 1.5878 = 0.000001 + 0.0003962 = 0.000397.
\end{align*}
The data changes at each rotation; the schedule does not. A rotation resumes the scheduler from the checkpoint, so epoch 21
continues from epoch 20's learning rate rather than jumping back to 0.0005.

\begin{sourcenote}[: no warm-up]
The F001 README once described ``one cosine schedule over 200 epochs (5 warm-up)''. The run's \code{config.json} sets
\code{warmup_epochs: 0}, and \code{FoundationTrainer._build_scheduler}
(\file{app/abstract_classes/foundation_trainer.py}) builds a plain \code{CosineAnnealingLR} that never reads that field. The
source wins: F001 has no warm-up, and its first epoch ran at the full 0.0005. The README has since been corrected (``no
warm-up: \texttt{warmup\_epochs} is 0 in the config'').
\end{sourcenote}

\coursefig[0.95\linewidth]{iw-43-rotation}{(a) The ten 30,000-patch blocks along F001's 200 epochs, under the one cosine
learning-rate curve that runs through all of them. (b) The loop in \texttt{rotate\_f001.sh}: every decision is computed
from \texttt{metrics.csv} and two marker files, so the script can be killed and restarted at any point.}{fig:pretrain:rotation}

\subsection{Early stopping is off}

F001 sets \code{early_stop_patience: null}, which turns early stopping off (the stop test begins with
\code{bool(cfg.early_stop_patience)}, which is false for \code{null}). The reason is the rotation. Within a block the model
sees the same 30,000 patches 20 times, so it can flatten out near the end of a block and improve again when fresh data
arrives. Suppose patience 10 and a best score at epoch 30:
\begin{center}
\begin{tabular}{rl}
\toprule
epochs 31--40 & no 1\% improvement; bad epochs climb 1, \dots, 10\\
epoch 40 & patience reached, $40 \ge 20$: \textbf{stop}\\
epoch 41 & block 2 would have arrived; never trained\\
\bottomrule
\end{tabular}
\end{center}
\code{best.pt} still tracks the lowest validation score, so the best epoch is kept either way.

\subsection{The driver: state read from the run's own records}

\file{research/experiments/indradhanu-se-pretrain/F001-final-pretrain/rotate_f001.sh} runs on the instance. The trainer
trains on whatever shard files sit in its local train folder, and \code{auto_resume} continues from \code{last.pt} with the
optimiser, scheduler and early-stop state. So a rotation is: stop, swap the folder's contents, start again. Each pass of the
driver's loop does five things:
\begin{enumerate}
  \item \textbf{Count the epochs done.} \code{metrics.csv} has a header plus one row per finished epoch, so
    $\text{done} = \text{lines} - 1$.
  \item \textbf{Pick the block.} $k = \lfloor \text{done} / 20 \rfloor$ (integer division) and
    $\text{stop\_at} = 20(k+1)$.
  \item \textbf{Rotate if needed.} A marker file \code{BLOCK} says which block the train folder holds. If it is not $k$, swap
    in the staged copy (only if a second marker says the stage holds block $k$ and all 200 files are there); otherwise
    download block $k$ now.
  \item \textbf{Prefetch} block $k + 1$ into a staging folder in the background.
  \item \textbf{Train}: start the trainer, poll every 30 seconds, and stop it once $\text{done} \ge \text{stop\_at}$.
\end{enumerate}

\begin{worked}[label={we:pretrain:driver}]{One pass of the driver after a crash}
The instance restarts after a crash and \code{metrics.csv} has 41 lines. Do each step:
\begin{align*}
  \text{done} &= 41 - 1 = 40\\
  k &= \lfloor 40 / 20 \rfloor = 2\\
  \text{stop\_at} &= 20 \times (2 + 1) = 60.
\end{align*}
\code{BLOCK} says 1, not 2, so the driver swaps in block 2 (shards $200 \times 2 = 400$ to $400 + 199 = 599$), starts
prefetching block 3 (shards 600--799), and resumes from \code{last.pt} (epoch 40): epochs 41 to 60, stopped at 60.
\end{worked}

\paragraph{Why stopping is safe.} The trainer saves the checkpoint and uploads it before it appends the epoch's row to
\code{metrics.csv}. So by the time the driver sees row 60, epoch 60's checkpoint is already safe. A state machine whose only
state is ``how many rows are in this file'' can be killed and restarted at any moment.

\paragraph{What it costs.} About 9.2 minutes per 30,000-patch epoch on the L40S at \$2.69 per hour:
\begin{align*}
  200 \times 9.2 \text{ min} &= 1{,}840 \text{ min}\\
  1{,}840 / 60 &= 30.7 \text{ hours}\\
  30.7 \times \$2.69 &= \$82.5,
\end{align*}
matching the README's pre-launch estimate of ``roughly 30 hours, \textasciitilde\$82''.\tbrk\tagM{research/experiments/indradhanu-se-pretrain/F001-final-pretrain/README.md}

\subsection{What the curves showed}

Two scores answer two questions. Validation L1 and SAM ask how well the model redraws SpectralEarth patches it never trained
on. The PRISMA bench asks how well it \emph{finds anomalies}: four labelled PRISMA scenes on PRISMA's own bands, scored by
mean ROC-AUC under the old \code{combined} score and, from epoch 52, under residual Mahalanobis (\cref{ch:residualmaha}).

\begin{figure}[htbp]
\centering
\begin{tikzpicture}
\begin{groupplot}[group style={group size=2 by 1, horizontal sep=1.9cm},
  width=0.47\linewidth, height=5.4cm, xmin=0, xmax=115, xlabel={epoch},
  tick label style={font=\scriptsize}, label style={font=\small}, legend style={font=\scriptsize}]
\nextgroupplot[ylabel={validation L1}, ymin=0.017, ymax=0.025, scaled y ticks=false,
  title={\small (a) reconstruction}, yticklabel style={/pgf/number format/fixed,/pgf/number format/precision=3}]
\addplot[spec, thick, mark=*] coordinates {(6,0.0241) (13,0.0219) (20,0.0202) (38,0.0192) (60,0.0185) (109,0.0185)};
\nextgroupplot[ylabel={mean ROC-AUC, PRISMA bench}, ymin=0.62, ymax=0.95, title={\small (b) detection},
  legend pos=north west, yticklabel style={/pgf/number format/fixed,/pgf/number format/precision=2}]
\addplot[loss, thick, mark=square*] coordinates {(7,0.664) (14,0.660) (20,0.734) (38,0.727) (60,0.734) (109,0.706)};
\addlegendentry{combined}
\addplot[spec, thick, mark=*] coordinates {(60,0.861) (109,0.862)};
\addlegendentry{residual Mahalanobis}
\addplot[maskgrey, dashed, thick, domain=0:115] {0.662};
\addlegendentry{v0.2.0, combined}
\end{groupplot}
\end{tikzpicture}

\begin{tikzpicture}
\begin{axis}[width=0.47\linewidth, height=4.6cm, xmin=0, xmax=115, xlabel={epoch}, ylabel={validation SAM (deg)},
  ymin=3.8, ymax=6.6, tick label style={font=\scriptsize}, label style={font=\small}, title={\small (c) spectral angle}]
\addplot[attn, thick, mark=triangle*] coordinates {(6,6.40) (13,5.53) (20,4.90) (38,4.47) (60,4.18) (109,4.16)};
\end{axis}
\end{tikzpicture}
\caption{F001's records drawn to scale from its README (validation rows at epochs 6, 13, 20, 38, 60, 109; bench at 7, 14,
20, 38, 60, 109). (a) Validation L1 falls at every checkpoint until epoch 60 and then flattens. (b) The bench under the
\texttt{combined} score jumps above production's 0.662 by epoch 20 and then stays flat; residual Mahalanobis, measured at
epochs 60 and 109, sits about 0.13 higher. (c) Validation SAM follows L1.}
\label{fig:pretrain:f001curves}
\end{figure}

Source for \cref{fig:pretrain:f001curves}:\tbrk\tagM{research/experiments/indradhanu-se-pretrain/F001-final-pretrain/README.md}.
From epoch 20 to 60 reconstruction improved:
\begin{align*}
  \text{L1:}&\quad 0.0202 - 0.0185 = 0.0017;\quad 0.0017 / 0.0202 = 0.084 \approx 8\% \text{ better}\\
  \text{SAM:}&\quad 4.90 - 4.18 = 0.72^\circ;\quad 0.72 / 4.90 = 0.147 \approx 15\% \text{ better.}
\end{align*}
Detection under \code{combined} did not move: 0.734, 0.727, 0.734. There are two ways to read a stall like that: the model
has stopped learning what matters, or the ruler cannot see the improvement. The same reconstructions scored with residual
Mahalanobis gave 0.861 at epoch 60, far above \code{combined}'s 0.734. The stall was the ruler: \code{combined} was dominated
by errors every pixel shares (detector stripes, haze, dark sea), so reducing the model's error on ordinary ground barely moved
it. \Cref{ch:residualmaha} derives why.

\coursefig[0.95\linewidth]{iw-43-curves}{The course's version of the same records, including the residual-Mahalanobis
measurement at epoch 52 (0.865) reported in the scoring comparison of 2026-10-15.}{fig:pretrain:coursecurves}

\begin{drift}
F001 was planned for 200 epochs and stopped at epoch 113 by decision, because the bench under residual Mahalanobis was flat:
0.861 at epoch 60, 0.862 at epoch 109.\tbrk\tagM{docs/proposals/fine-tuning-and-segmentation.md} Fine-tuning starts from the
epoch-109 checkpoint, \code{F001_ep109.pt}. Under residual Mahalanobis, production's v0.2.0 scored 0.854, only 0.011 behind F001's
epoch 52 (0.865), and caught more objects (403 against 383 of 909) (IW43, 2026-10-15). One scene can swing a mean: set1 scored
0.430 at epoch 14 and 0.702 at epoch 20 under \code{combined}, a swing of 0.27 from one block of training. Four scenes cannot
give a reliable mean when one of them moves that much, which is why every later gate compares runs on the same scenes with the
same score (\cref{ch:bootstrap}).
\end{drift}

\section{Three sensors at their own bands}\label{sec:pretrain:native}

\subsection{Why new data at all}

F001 has only seen SpectralEarth: EnMAP data on one fixed list of 202 bands. The product scores three hyperspectral sensors,
PRISMA, EnMAP and AVIRIS-NG (\cref{ch:sensors}), each on its own band list and at its own pixel size. Fine-tuning
(\cref{ch:finetune}) trains on them directly, for three reasons given in
\file{docs/proposals/fine-tuning-and-segmentation.md}: real sensors; the wavelengths SpectralEarth dropped (including the
hydrocarbon band near 1730\,nm); and EnMAP's cloud and water labels for a segmentation head (\cref{ch:seghead}).

\begin{definition}[Native and resampled]\index{native band list}
A \emph{native} band list is the sensor's own list of band centres and its own pixel grid, untouched. A \emph{resampled} list
holds values interpolated onto some other list, such as the 165-band grid of \cref{ch:sensors} or SpectralEarth's 202.
\end{definition}

All three sets share SpectralEarth's sample layout: \code{pixels.npy} (or compressed \code{pixels.i16dz}), int16
\shape{(bands, 128, 128)}, multiplied by 0.0001 for reflectance, $-32768$ for nodata; \code{wavelengths.npy}, float32
\shape{(bands,)}, the centre of each band in the same order; \code{meta.json} with scene facts; and for EnMAP the provider's
quality layers \code{label_*.npy}. A tile is kept only if at least 10\% of its pixels are valid.

\coursefig[0.95\linewidth]{iw-44-datasets}{(a) The three shard sets as they stood on 2026-10-16 (colours identify the sensor
only). (b) The one rule that lets mixed band lists share a trainer: match by wavelength within 0.5\,nm; a target band with no
match becomes nodata, a source band with no target is dropped.}{fig:pretrain:datasets}

\subsection{EnMAP: 224 bands, labels and a defect map}

210 EnMAP L2A scenes became 14,542 training and 950 test tiles, each with all 224 bands at the scene's own
wavelengths.\tbrk\tagM{docs/lld/enmap-shards.md}

\paragraph{A renumbered band.} Five early-2023 scenes (processor version V010502) number their shortwave bands one off: there
band 142 is 1530.29\,nm, elsewhere 1518.87\,nm. Reading ``band 142'' by position would silently mix two wavelengths. So every
sample carries its wavelengths, and everything downstream matches \emph{by wavelength, never by band number}. Across all scenes
the stats list 227 distinct wavelengths: 224 per scene plus the renumbered scenes' shifted ones.

\paragraph{The split by datatake.} A \emph{datatake}\index{datatake} is one continuous recording; consecutive scenes from it
border each other. Whole datatakes go to one side of the split: 6 of the 66 went to test (13 scenes), with cloudy and clear
ones represented. A random split by scene would put ground a few pixels from a training scene into the test set.

\paragraph{The pixel mask.} EnMAP's \code{QL_PIXELMASK} flags individual voxels (one band of one pixel); the tiler turns each
flagged voxel into nodata. A probe of three scenes found bands 131 to 135 (about 1343--1390\,nm, the L2A water-vapour blank)
flagged everywhere, and about 200 other bands with scattered flags whose share is the same in all three scenes to 0.01\%
(band 187: 5.31\%). A share that does not change with the scene follows the \emph{detector}, not the ground: a fixed defect map.
Net effect: 85--95\% of pixels lose 2--4 of their 224 bands. That is harmless: a dead band drops out of the wavelength stem
exactly (\cref{ch:crossstem}) and out of the loss too. A pixel with 221 of 224 bands is a slightly thinner spectrum, not a hole.

\paragraph{Labels.} EnMAP's L2A product ships quality layers: cloud, cirrus (levels 0--3), haze, cloud shadow, snow, and a class
map. They are stored exactly as delivered and never applied to the pixels: a cloudy tile stays in, because cloudy tiles are what
a segmentation head learns from. \Cref{ch:seghead} audits them.

\subsection{PRISMA: reading the file without the product pipeline}

223 PRISMA L2D scenes became 14,802 training and 1,570 test tiles.\tbrk\tagM{docs/lld/prisma-shards.md} The product reads PRISMA
through a builder that resamples onto the 165-band grid; the shards must keep the sensor's bands, so \code{read_scene}
(\file{app/utils/prisma/scene.py}) opens the HDF5 file directly and applies three rules:
\begin{align*}
  \text{reflectance} &= \frac{\text{DN}}{65535}\,\bigl(\text{Max}_{\text{det}} - \text{Min}_{\text{det}}\bigr) + \text{Min}_{\text{det}}\\
  \text{nodata} &\iff \text{DN} = 0 \text{ or the detector's error flag} \ne 0\\
  \text{dropped} &\iff \text{the band's flag is 0 (5 of 239 bands).}
\end{align*}
DN (digital number\index{digital number}) is the stored integer, 0 to 65,535; the detector \emph{det} is VNIR or SWIR, each with its own scale
factors in the file. So $239 - 5 = 234$ bands remain.

\begin{worked}[label={we:pretrain:dn}]{From a PRISMA digital number to a stored value}
Take DN $= 13{,}107$ with illustrative factors \tagA\ Min $= 0.0$, Max $= 1.0$.
\begin{align*}
  13{,}107 / 65{,}535 &= 0.2\\
  \text{reflectance} &= 0.2 \times (1.0 - 0.0) + 0.0 = 0.2\\
  \text{stored} &= \operatorname{round}(0.2 / 0.0001) = \operatorname{round}(2000) = 2{,}000 \text{ (int16).}
\end{align*}
The LLD's check: on set1 at 860\,nm the result differs from the product builder's own read by at most 0.00005 reflectance, which
is int16 rounding: one stored unit is 0.0001, so rounding moves a value by at most half of that.
\end{worked}

\paragraph{Odd band order, one list everywhere.} The file lists the VNIR bands from about 980\,nm \emph{downwards}, then the SWIR
bands upwards, and the detectors overlap between about 920 and 980\,nm. Nothing is sorted or de-duplicated: the stem takes bands
in any order because each carries its wavelength. All 234 centres were identical to 0.000\,nm across every scene checked, so
PRISMA batches need no alignment.

\paragraph{234 or 185?} The shards hold 234 PRISMA bands; the bench reads 185, PRISMA's own bands after the product's default band
filter removes noisy and overlapping ones. Both are native; they are not the same list.

\paragraph{The split by site, single-link.} PRISMA revisits places, so \code{app/utils/prisma/split.py} groups scenes into sites:
scenes whose centres lie within 15\,km of each other, \emph{chained}. Single-link chaining\index{single-link clustering} means two scenes are in one site if a
chain of steps, each at most 15\,km, joins them:
\begin{center}
\begin{tabular}{lcl}
\toprule
pair & distance & \\
\midrule
P--Q & 10\,km & $\le 15$: same site\\
Q--R & 12\,km & $\le 15$: same site\\
P--R & 21\,km & $> 15$, but P and R are already joined through Q\\
\bottomrule
\end{tabular}
\end{center}
Result: one site $\{P, Q, R\}$. Any site within 30\,km of a benchmark scene is excluded from both sides (on the real run, none
was). Of the rest, 10\% of sites went to test: 21 scenes against 202 for training. PRISMA L2D has no cloud or water mask, so PRISMA
batches teach reconstruction only.

\subsection{AVIRIS-NG: an airborne sensor in four product versions}

AVIRIS-NG flies on an aircraft. The 79 flight lines of the ISRO--NASA India campaign became 22,377 training and 2,270 test
tiles.\tbrk\tagM{docs/lld/aviris-shards.md} Each line has 425 bands, of which the header's bad-band list marks 53 (the water-vapour
windows): $425 - 53 = 372$ remain. The campaign ships four product versions that place the 372 centres up to 0.42\,nm apart, so
across all lines the stats list 1,116 distinct wavelengths in 372 tight groups ($3 \times 372 = 1{,}116$: the 372 bands appear at
up to three positions). The fine-tuning band list \code{aviris_ng_372.json} takes the \emph{median} of each group; every variant is
within 0.5\,nm of it.

\paragraph{Pixel size.} 3.2 to 8.2\,m, depending on the line, deliberately not resampled to 30\,m. A tile is always 128 pixels
across, so the ground it covers changes:
\begin{align*}
  \text{EnMAP, PRISMA:}&\quad 128 \times 30 \text{ m} = 3{,}840 \text{ m}\\
  \text{AVIRIS-NG:}&\quad 128 \times 3.2 \text{ m} = 409.6 \approx 410 \text{ m}, \qquad 128 \times 8.2 \text{ m} = 1{,}049.6 \approx 1{,}050 \text{ m.}
\end{align*}
An AVIRIS-NG tile shows a field; a satellite tile shows a district.

\subsection{Making it fit: band-delta compression}

Neighbouring bands of one pixel are similar, so storing each band as its \emph{difference from the previous band} leaves small
numbers, which a general-purpose compressor (zstd) packs far better (\file{app/utils/shard_codec.py}). The format is
\code{.i16dz}.\index{band-delta compression}

\begin{worked}[label={we:pretrain:delta}]{Delta-encoding one pixel}
Four bands of one pixel, 1000, 1012, 1020, 1015. The encoder subtracts a zero band in front, so the first value is kept:
\begin{align*}
  \text{deltas:}\quad & 1000 - 0 = 1000,\quad 1012 - 1000 = 12,\\
                      & 1020 - 1012 = 8,\quad 1015 - 1020 = -5\\
  \text{decode (running sum):}\quad & 1000;\quad 1000 + 12 = 1012;\\
                      & 1012 + 8 = 1020;\quad 1020 - 5 = 1015.
\end{align*}
\end{worked}

\paragraph{Wrap-around is exact.} The differences are taken in int16, which holds $-32{,}768$ to $32{,}767$. Next to nodata a
difference overflows, and the codec lets it wrap. Wrapping is arithmetic \emph{modulo}\index{modular arithmetic} $2^{16} = 65{,}536$: a result is replaced
by the unique number in the int16 range that differs from it by a multiple of 65,536. Encoding and decoding both work modulo
65,536, and $(a - b) + b \equiv a \pmod{65{,}536}$, so the round trip returns $a$ whenever $a$ itself is in range:
\begin{align*}
  \text{encode:}\quad & 1000 - (-32{,}768) = 33{,}768 \quad\text{(too big for int16)}\\
  & 33{,}768 - 65{,}536 = -31{,}768 \quad\text{(stored, wrapped)}\\
  \text{decode:}\quad & -32{,}768 + (-31{,}768) = -64{,}536 \quad\text{(too small)}\\
  & -64{,}536 + 65{,}536 = 1000 \quad\text{(exact).}
\end{align*}

\paragraph{The size arithmetic.} AVIRIS-NG tiles were written compressed from the start; EnMAP and PRISMA were re-encoded on
2026-10-16, about 2.1 times smaller.\tbrk\tagM{research/experiments/indradhanu-finetune/README.md} From the sizes in the three shard LLDs
(GiB, train $+$ test):
\begin{align*}
  \text{EnMAP raw:}&\quad 100.9 + 6.6 = 107.5; \quad\text{re-encoded } 107.5 / 2.1 = 51.2\\
  \text{PRISMA raw:}&\quad 105.8 + 11.2 = 117.0; \quad\text{re-encoded } 117.0 / 2.1 = 55.7\\
  \text{AVIRIS-NG:}&\quad 116.0 + 11.2 = 127.2 \quad\text{(already compressed)}\\
  \text{total raw:}&\quad 107.5 + 117.0 + 127.2 = 351.7\\
  \text{total re-encoded:}&\quad 51.2 + 55.7 + 127.2 = 234.1.
\end{align*}
351.7\,GiB would leave no working room on a 412\,GB disk ($351.7 \times 1.0737 = 377.6$\,GB, since 1\,GiB $= 2^{30}$ bytes $=
1.0737$\,GB); 234.1\,GiB ($251.4$\,GB) fits, with room to spare.

\begin{sourcenote}[: two compression ratios]
SpectralEarth's shards compress 1.89--1.91 times (\file{docs/lld/spectralearth.md}; the docstring of
\file{app/utils/shard_codec.py} says 1.91x). The native-v1 EnMAP and PRISMA sets compress about 2.1 times
(\file{research/experiments/indradhanu-finetune/README.md}). Comments in \file{recompress_shards.py} and \file{recompress_ft.sh}
still quote SpectralEarth's \textasciitilde 1.9x for the fine-tuning sets. Both numbers are right for their own data; quote the one
that belongs to the shards you are sizing.
\end{sourcenote}

\subsection{One batch, one band list}\label{sec:pretrain:onelist}

The reduced stem takes \emph{one wavelength list per call}: a single \shape{(N,)} vector, not one per sample. A batch is a stack
of samples, so every sample in it must carry the same bands in the same order. Two rules make that true:
\begin{enumerate}
  \item \textbf{One sensor per batch} (relaxed in \cref{ch:finetune} to one band list per \emph{stem pass}).
  \item \textbf{Within a sensor, align by wavelength} (\file{app/utils/band_alignment.py}). Each target band takes the nearest
    source band within 0.5\,nm; a target band with no match becomes nodata; a source band with no target is dropped. Nothing is
    interpolated.
\end{enumerate}
The tolerance 0.5\,nm is chosen so it ``only absorbs rounding in the metadata and never pairs two neighbouring bands'': EnMAP's
bands are 6--12\,nm apart, AVIRIS-NG's about 5\,nm, so two real neighbours are always more than ten tolerances apart.

\begin{worked}[label={we:pretrain:align}]{Aligning a sample onto a band list}
Target list 500, 510, 520\,nm; a sample with bands at 500.2, 520.1 and 530.0\,nm.
\begin{center}
\begin{tabular}{lllll}
\toprule
target & nearest source & distance & test & result\\
\midrule
500 & 500.2 & $|500 - 500.2| = 0.2$ & $\le 0.5$ & take it\\
510 & 500.2 & $|510 - 500.2| = 9.8$ & $> 0.5$ & nodata\\
520 & 520.1 & $|520 - 520.1| = 0.1$ & $\le 0.5$ & take it\\
\bottomrule
\end{tabular}
\end{center}
Source 530.0 is near no target and is dropped. The aligned sample has three bands, one of them dead; the stem skips the dead one
exactly.
\end{worked}

Per sensor: EnMAP aligns onto \code{enmap_224.json}, the 224 centres of the standard scenes (the renumbered scenes' shifted bands
are dropped, and standard bands they lack become nodata); PRISMA needs nothing; AVIRIS-NG aligns onto the medians in
\code{aviris_ng_372.json}, and every variant pairs.

\begin{confusions}
\confusion{Should training use the full SAM weight from epoch 1 too?}{No. The ramp exists for the gradients (\cref{ch:losses}):
early SAM gradients are large and noisy. The fix changes only the number used to compare epochs, not what the model learns from.}
\confusion{Isn't a minimum of 20 epochs already protection against the trap?}{It delays the stop until the ramp has finished, but
epoch 1's 0.055 is still the best on record at epoch 20, so the stop fires there and \texttt{best.pt} is epoch 1.}
\confusion{So the wavelength-aware stem was a mistake?}{No. A fixed $1 \times 1$ compressor cannot even be \emph{asked} to read a
different list. What was missing was training that exercised the ability. Band-list augmentation provides it.}
\confusion{Isn't band-list augmentation the ``hiding bands'' trap of an earlier chapter?}{It hides bands from the \emph{input}, but
the loss is taken only on pixels the token mask hid completely, and a hidden pixel's own bands are never an input. A pixel never
learns to fill in its own missing bands, so an anomaly cannot rebuild itself.}
\confusion{Does rotating the data restart the training?}{No. Weights, optimiser state and learning rate all continue from
\texttt{last.pt}. Only the files in the train folder change.}
\confusion{Why 20 epochs per block, not 1?}{A block is one epoch's worth of patches, so one epoch per block would mean a 198\,GB
download every 9 minutes. Twenty epochs per block gives the prefetch about three hours and keeps each patch's repeat count equal to
the E-runs'.}
\confusion{Does nodata mean zero?}{No. Nodata is $-32{,}768$ in the file and becomes an invalid voxel in the trainer, which the stem
skips exactly. A zero would be read as ``black in this band''.}
\confusion{PRISMA has 234 bands or 185?}{Both. The shards hold all 234 real bands; the bench and the product read 185 after the
default band filter.}
\end{confusions}

\begin{checkyourself}
\item\label{ex:pretrain:1} Using the table of \cref{we:pretrain:ramp}, a later epoch has L1 $= 0.050$ and SAM $= 0.18$ rad at
  $\lambda = 0.05$. Compute its score with the ramp and at the full weight, and use the decomposition
  $\Delta L1 + \lambda_e\Delta\SAM + \Delta\lambda\,\SAM_{e+1}$ to explain the change from epoch 1 under the ramp.
\item\label{ex:pretrain:2} Epoch 1 scores 0.155 at the full weight. What must a later epoch score to count as an improvement under
  $\delta = 0.01$? If every later epoch scores 0.1540, on which epoch does the run stop (patience 10, minimum 20)?
\item\label{ex:pretrain:3} A sensor C samples the visible three times (500, 505, 510\,nm, value 2) and the infrared once (2000\,nm,
  value $-1$). Compute the plain $1/N$ skip output, then the width-weighted output with idealised widths of 5\,nm per visible band
  and 10\,nm for the infrared band. Explain why the width-weighted answer is not 0.5.
\item\label{ex:pretrain:4} After a restart \code{metrics.csv} has 128 lines. Work out done, $k$, stop\_at and the shards of the block,
  and compute the learning rate at $t = 100$ from \eqref{eq:pretrain:cosine}.
\item\label{ex:pretrain:5} Delta-encode the bands 2000, 2010, 1990, 1995, then decode them. Then encode the pair (nodata $-32{,}768$,
  then 30{,}000) with int16 wrap-around and decode it again, showing each step.
\item\label{ex:pretrain:6} Align a sample with bands 500.3, 509.6, 525.0\,nm onto the list 500, 510, 520\,nm. Which bands are kept,
  which targets become nodata, which source bands are dropped?
\end{checkyourself}

\nextchapter{ch:seghead}
\endgroup
